\input{../tex-inputs/latex-leseansicht-vorspann}

\section[Gustav Schwarzkopf an Arthur Schnitzler, 14. 4. 1906]{L04269 Gustav Schwarzkopf an Arthur Schnitzler, 14. 4. 1906}
\nopagebreak\mylabel{L04269v}
\rehead{ }\normalsize\beginnumbering\briefempfaengerindex{Schnitzler, Arthur@\textsc{Schnitzler, Arthur}!zzzSchwarzkopf, Gustav@\emph{von Gustav Schwarzkopf}!1906-04-141@{14. 4. 1906}|(be}
\toendnotes[C]{\smallbreak\pagebreak[2]}
\correspDesc{Versand  durch Gustav Schwarzkopf am 14. 4. 1906 in Wien
\newline{}Übermittlung  am 14. 4. 1906 in Wien
\newline{}Erhalt  durch Arthur Schnitzler am 14. 4. 1906 in Wien}\toendnotes[C]{\smallbreak}
\Standort{CUL, Schnitzler, B 96a.}
\physDesc{Postkarte, 569 Zeichen
\newline{}Handschrift: schwarze Tinte, deutsche Kurrent
\newline{}Versand: 1) Stempel: »\nobreak{}\oindex{Wien@\textbf{Wien}, \emph{Verwaltungsgebiet}|pwk}{[}Wie{]}n, 16 IV 06, 12–V\nobreak{}«.   2) Stempel: »\nobreak{}\oindex{XVIII., Währing@\textbf{XVIII., Währing}, \emph{Verwaltungsgebiet}|pwk}18/\textsubscript{1} Wien 111, 16 IV \textcolor{gray}{0}6, XII\textsuperscript{40}\nobreak{}«. }\toendnotes[C]{\smallbreak}\pstart{}{\pb}I. Tiefer Graben\pend{}\pstart{}23. II. St. Th. 6.\pend{}{\bigskip}\pstart{}Herrn \textsc{D\textsuperscript{r} Arthur
                     Schnitzler}\pend{}\pstart{}\textsc{Wien\oindex{Wien@\textbf{Wien}, \emph{Verwaltungsgebiet}|pw}}\pend{}\pstart{}XVIII Spöttelgaſſe 7\oindex{Wien@\textbf{Wien}!XVIII., Währing@\textbf{XVIII., Währing}!Edmund-Weiß-Gasse 7@\textbf{Edmund-Weiß-Gasse 7}, \emph{Wohngebäude}|pw}.\pend{}{\bigskip}\vspace{1em}
\pstart
           \raggedleft{}{\pb}14. IV. 06\pend
           \vspace{0.5em}
\pstart
           Lieber Arthur, wir werden Sie \label{K_L04269-1v}\edtext{morgen, (So{\geminationn}tag}{\lemma{\textnormal{\emph{morgen, (Sonntag}}}\Cendnote{\textnormal{Vgl. A. S.: \emph{Tagebuch}, 15. 4. 1906.
               }}}\label{K_L04269-1}) bei
               gutem Wetter um ½ 10 Uhr abholen, bitten Sie aber es nicht übel zu
               nehmen, we{\geminationn} wir Ihre, freundliche Einladung, zum
               Mittageſſen zu ko{\geminationm}en, nicht annehmen. Nach einem
               längeren Spaziergang bin ich als Gaſt abſolut nicht mehr zu verwenden und auch 
                  Max\pwindex{Schwarzkopf, Max 12.\,6.\,1857 Wien – 14.\,4.\,1928 ebd.@\textsc{Schwarzkopf, Max} (12.\,6.\,1857 Wien – 14.\,4.\,1928 ebd.), \emph{Rechtsanwalt}|pw}
                macht es{ }ſich nach dem Essen gerne ein wenig bequem. Bei{ }ſchlechtem Wetter aber
               will ich gar \uline{nicht ausgehen}, weil meine
                  \textcolor{gray}{Unterh} u ſ. \textcolor{gray}{w}. Es wird aber hübſch{ }ſein.\pend
           
\pstart
           Herzlichſt grüßend{\\[\baselineskip]}Ihr{\\[\baselineskip]}\spacefill\mbox{Gustav.}\pend
           \leftskip=0em{}\selectlanguage{ngerman}\endnumbering\briefempfaengerindex{Schnitzler, Arthur@\textsc{Schnitzler, Arthur}!zzzSchwarzkopf, Gustav@\emph{von Gustav Schwarzkopf}!1906-04-141@{14. 4. 1906}|)be}\mylabel{L04269h}
\begin{anhang}
\end{anhang}\newcommand{\dateiname}{L04269}\newcommand{\titel}{Gustav Schwarzkopf an Arthur Schnitzler, 14. 4. 1906}\newcommand{\editorInnen}{Herausgegeben von Jahnke, SelmaMüller, Martin Anton}\input{../tex-inputs/latex-leseansicht-abspann}
